\documentclass[11pt,letterpaper]{article}
\usepackage[technical-reference]{cornell-notes}
\usepackage{needspace}

\title{Clause 14: Exception Handling - Cornell Notes}
\author{}
\date{}
\setDocTitle{Clause 14: Exception Handling}
\setDocSubtitle{C++ 2024 Cornell Notes}
\setDocOwner{C++ 2024 Cornell Notes}
\setDocAuthor{}
\setDocDate{}
\setCornellCollection{C++ 2024 Cornell Notes}
\setCornellUnitType{Clause}
\setCornellUnitNumber{14}
\setCornellUnitTitle{Exception Handling}
\setCornellCompanionLabel{C++ Technical Reference}
\hypersetup{pdftitle={Clause 14: Exception Handling - Cornell Notes},pdfsubject={C++ technical study notes},pdfauthor={}}

\begin{document}
\maketitle
\begin{CornellReferenceOrientationBox}
\textbf{Purpose.} Defines throwing, handler matching, stack unwinding, constructor and destructor behavior, exception specifications, termination, and uncaught-exception observation.
\par\textbf{Role.} Specifies non-local error transfer and cleanup semantics.
\par\textbf{Principal dependencies.} Uses expressions, object lifetime, functions, classes, and type matching.
\par\textbf{Primary audience.} programmers, runtime implementers, and exception-safety reviewers.
\par\textbf{Major subject groups.} Preamble; Throwing an exception; Constructors and destructors; Handling an exception; Exception specifications; Special functions.
\end{CornellReferenceOrientationBox}
\section{Learning Objectives}
\begin{itemize}[label=\textcolor{CornellReferenceTeal}{$\blacktriangleright$}]
\item Define the governing ideas of Preamble and state the consequences for a program or implementation.
\item Identify the governing ideas of Throwing an exception and state the consequences for a program or implementation.
\item Distinguish the governing ideas of Constructors and destructors and state the consequences for a program or implementation.
\item Explain the governing ideas of Handling an exception and state the consequences for a program or implementation.
\item Trace the governing ideas of Exception specifications and state the consequences for a program or implementation.
\item Apply the governing ideas of Special functions and state the consequences for a program or implementation.
\item Analyze the governing ideas of General and state the consequences for a program or implementation.
\item Evaluate the governing ideas of The std::terminate function and state the consequences for a program or implementation.
\item Diagnose the governing ideas of The std::uncaught\_\allowbreak{}exceptions function and state the consequences for a program or implementation.
\item Compare the governing ideas of exception object and state the consequences for a program or implementation.
\end{itemize}
\section{Normative Structure Map}
\small\begin{longtable}{@{}>{\RaggedRight\arraybackslash}p{0.13\textwidth}>{\RaggedRight\arraybackslash}p{0.27\textwidth}>{\RaggedRight\arraybackslash}p{0.19\textwidth}>{\RaggedRight\arraybackslash}p{0.31\textwidth}@{}}
\toprule\textbf{Subclause} & \textbf{Subject} & \textbf{Rule category} & \textbf{Key dependencies / entities}\\\midrule\endfirsthead
\toprule\textbf{Subclause} & \textbf{Subject} & \textbf{Rule category} & \textbf{Key dependencies / entities}\\\midrule\endhead
\bottomrule\endfoot
14.1 & Preamble & core-language semantics & Uses expressions, object lifetime, functions, classes, and type matching.\\
14.2 & Throwing an exception & error behavior & Uses expressions, object lifetime, functions, classes, and type matching.\\
14.3 & Constructors and destructors & object contract & Uses expressions, object lifetime, functions, classes, and type matching.\\
14.4 & Handling an exception & error behavior & Uses expressions, object lifetime, functions, classes, and type matching.\\
14.5 & Exception specifications & error behavior & Uses expressions, object lifetime, functions, classes, and type matching.\\
14.6 & Special functions & operation semantics & General, The std::terminate function, The std::uncaught\_\allowbreak{}exceptions function\\
\end{longtable}\normalsize
\begin{CornellReferenceNormativeBox}A heading maps the clause; it does not replace the detailed conditions. For each operation, read requirements, effects, result, exceptions, complexity, remarks, and notes with their stated normative force.\end{CornellReferenceNormativeBox}
\subsection*{Rule Classification Guide}
\small\begin{longtable}{@{}>{\RaggedRight\arraybackslash}p{0.25\textwidth}>{\RaggedRight\arraybackslash}p{0.67\textwidth}@{}}\toprule\textbf{Label or category} & \textbf{How to read it}\\\midrule\endfirsthead\toprule\textbf{Label or category} & \textbf{How to read it (continued)}\\\midrule\endhead\bottomrule\endfoot
Well-formed / ill-formed & Formation is governed by syntax and semantic rules; an ill-formed program does not become well-formed because one implementation accepts it.\\
Diagnosable rule & A violation requires at least one diagnostic unless the wording places the case outside the diagnosable set.\\
No diagnostic required & The program can be ill-formed while the implementation has no diagnostic obligation.\\
Undefined behavior & No execution requirements are imposed; translation success does not establish safety or portability.\\
Unspecified behavior & One of the permitted outcomes may be chosen without documentation.\\
Implementation-defined behavior & The implementation chooses and documents the behavior.\\
Conditionally supported & The implementation may omit support but documents unsupported constructs.\\
\end{longtable}\normalsize
\section{Cornell Cue-and-Notes}
\Needspace{8\baselineskip}
\subsection*{14.1 Preamble}
\begin{longtable}{@{}>{\RaggedRight\arraybackslash\columncolor{CornellReferenceCueGray}}p{0.28\textwidth}>{\RaggedRight\arraybackslash}p{0.64\textwidth}@{}}
\rowcolor{CornellReferenceNavy}\color{white}\textbf{Cue / recall prompt} & \color{white}\textbf{Notes}\\\endfirsthead
\rowcolor{CornellReferenceNavy}\color{white}\textbf{Cue / recall prompt} & \color{white}\textbf{Notes (continued)}\\\endhead
\arrayrulecolor{CornellReferenceLineGray}
\textbf{What rule applies to Preamble?}\\[-1mm]{\scriptsize\CornellClauseRef{14.1}} & Frames the terminology, applicability, and relationships needed to interpret Preamble; detailed rules remain in the following subclauses.\\\hline
\end{longtable}
\begin{CornellReferenceSummaryBox}Preamble supplies the core-language semantics portion of this clause. The key review move is to connect its local wording to the clause-wide model, then classify each condition by normative force before relying on it.\end{CornellReferenceSummaryBox}
\Needspace{8\baselineskip}
\subsection*{14.2 Throwing an exception}
\begin{longtable}{@{}>{\RaggedRight\arraybackslash\columncolor{CornellReferenceCueGray}}p{0.28\textwidth}>{\RaggedRight\arraybackslash}p{0.64\textwidth}@{}}
\rowcolor{CornellReferenceNavy}\color{white}\textbf{Cue / recall prompt} & \color{white}\textbf{Notes}\\\endfirsthead
\rowcolor{CornellReferenceNavy}\color{white}\textbf{Cue / recall prompt} & \color{white}\textbf{Notes (continued)}\\\endhead
\arrayrulecolor{CornellReferenceLineGray}
\textbf{What rule applies to Throwing an exception?}\\[-1mm]{\scriptsize\CornellClauseRef{14.2}} & Defines the failure model for Throwing an exception, including how an error is represented, reported, propagated, or converted into a diagnostic consequence.\\\hline
\end{longtable}
\begin{CornellReferenceSummaryBox}Throwing an exception supplies the error behavior portion of this clause. The key review move is to connect its local wording to the clause-wide model, then classify each condition by normative force before relying on it.\end{CornellReferenceSummaryBox}
\Needspace{8\baselineskip}
\subsection*{14.3 Constructors and destructors}
\begin{longtable}{@{}>{\RaggedRight\arraybackslash\columncolor{CornellReferenceCueGray}}p{0.28\textwidth}>{\RaggedRight\arraybackslash}p{0.64\textwidth}@{}}
\rowcolor{CornellReferenceNavy}\color{white}\textbf{Cue / recall prompt} & \color{white}\textbf{Notes}\\\endfirsthead
\rowcolor{CornellReferenceNavy}\color{white}\textbf{Cue / recall prompt} & \color{white}\textbf{Notes (continued)}\\\endhead
\arrayrulecolor{CornellReferenceLineGray}
\textbf{How is Constructors and destructors initialized?}\\[-1mm]{\scriptsize\CornellClauseRef{14.3}} & Specifies how Constructors and destructors establishes object state, including participation constraints, initialization effects, and exceptional exits.\\\hline
\end{longtable}
\begin{CornellReferenceSummaryBox}Constructors and destructors supplies the object contract portion of this clause. The key review move is to connect its local wording to the clause-wide model, then classify each condition by normative force before relying on it.\end{CornellReferenceSummaryBox}
\Needspace{8\baselineskip}
\subsection*{14.4 Handling an exception}
\begin{longtable}{@{}>{\RaggedRight\arraybackslash\columncolor{CornellReferenceCueGray}}p{0.28\textwidth}>{\RaggedRight\arraybackslash}p{0.64\textwidth}@{}}
\rowcolor{CornellReferenceNavy}\color{white}\textbf{Cue / recall prompt} & \color{white}\textbf{Notes}\\\endfirsthead
\rowcolor{CornellReferenceNavy}\color{white}\textbf{Cue / recall prompt} & \color{white}\textbf{Notes (continued)}\\\endhead
\arrayrulecolor{CornellReferenceLineGray}
\textbf{What rule applies to Handling an exception?}\\[-1mm]{\scriptsize\CornellClauseRef{14.4}} & Defines the failure model for Handling an exception, including how an error is represented, reported, propagated, or converted into a diagnostic consequence.\\\hline
\end{longtable}
\begin{CornellReferenceSummaryBox}Handling an exception supplies the error behavior portion of this clause. The key review move is to connect its local wording to the clause-wide model, then classify each condition by normative force before relying on it.\end{CornellReferenceSummaryBox}
\Needspace{8\baselineskip}
\subsection*{14.5 Exception specifications}
\begin{longtable}{@{}>{\RaggedRight\arraybackslash\columncolor{CornellReferenceCueGray}}p{0.28\textwidth}>{\RaggedRight\arraybackslash}p{0.64\textwidth}@{}}
\rowcolor{CornellReferenceNavy}\color{white}\textbf{Cue / recall prompt} & \color{white}\textbf{Notes}\\\endfirsthead
\rowcolor{CornellReferenceNavy}\color{white}\textbf{Cue / recall prompt} & \color{white}\textbf{Notes (continued)}\\\endhead
\arrayrulecolor{CornellReferenceLineGray}
\textbf{What rule applies to Exception specifications?}\\[-1mm]{\scriptsize\CornellClauseRef{14.5}} & Defines the failure model for Exception specifications, including how an error is represented, reported, propagated, or converted into a diagnostic consequence.\\\hline
\end{longtable}
\begin{CornellReferenceSummaryBox}Exception specifications supplies the error behavior portion of this clause. The key review move is to connect its local wording to the clause-wide model, then classify each condition by normative force before relying on it.\end{CornellReferenceSummaryBox}
\Needspace{8\baselineskip}
\subsection*{14.6 Special functions}
\begin{longtable}{@{}>{\RaggedRight\arraybackslash\columncolor{CornellReferenceCueGray}}p{0.28\textwidth}>{\RaggedRight\arraybackslash}p{0.64\textwidth}@{}}
\rowcolor{CornellReferenceNavy}\color{white}\textbf{Cue / recall prompt} & \color{white}\textbf{Notes}\\\endfirsthead
\rowcolor{CornellReferenceNavy}\color{white}\textbf{Cue / recall prompt} & \color{white}\textbf{Notes (continued)}\\\endhead
\arrayrulecolor{CornellReferenceLineGray}
\textbf{What rule applies to General?}\\[-1mm]{\scriptsize\CornellClauseRef{14.6.1}} & Frames the terminology, applicability, and relationships needed to interpret General; detailed rules remain in the following subclauses.\\\hline
\textbf{What rule applies to The std::terminate function?}\\[-1mm]{\scriptsize\CornellClauseRef{14.6.2}} & Defines the core-language rules for The std::terminate function, including formation, semantic effect, interactions with type and lookup, and any ill-formed or implementation-dependent cases.\\\hline
\textbf{What rule applies to The std::uncaught\_\allowbreak{}exceptions function?}\\[-1mm]{\scriptsize\CornellClauseRef{14.6.3}} & Defines the failure model for The std::uncaught\_\allowbreak{}exceptions function, including how an error is represented, reported, propagated, or converted into a diagnostic consequence.\\\hline
\end{longtable}
\begin{CornellReferenceSummaryBox}Special functions supplies the operation semantics portion of this clause. The key review move is to connect its local wording to the clause-wide model, then classify each condition by normative force before relying on it.\end{CornellReferenceSummaryBox}
\section{Language-Lawyer and Library-Usage Pitfalls}
\begin{CornellReferencePitfallBox}\begin{itemize}[label=\textcolor{CornellReferenceAmber}{$\blacktriangleright$}]
\item Throwing from cleanup without considering an already active exception.
\item Assuming a destructor always runs for a not-fully-constructed object.
\item Confusing handler selection with overload resolution.
\item Treating a non-throwing specification as a promise the runtime can ignore.
\item Treating a note, example, or recommended practice as though it imposed a mandatory program rule.
\end{itemize}\end{CornellReferencePitfallBox}
\section{Programmer and Implementation Checklist}
\begin{CornellReferenceChecklistBox}\begin{itemize}[label=$\square$]
\item Identify the exception object's type and initialization.
\item Determine the first matching handler in the permitted search.
\item List objects destroyed during unwinding.
\item Audit constructors and destructors for propagation behavior.
\item Check potentially-throwing status and termination cases.
\item For \CornellClauseRef{14.1}, verify preamble against the precise rule category and all cited dependencies.
\item For \CornellClauseRef{14.2}, verify throwing an exception against the precise rule category and all cited dependencies.
\item For \CornellClauseRef{14.3}, verify constructors and destructors against the precise rule category and all cited dependencies.
\item For \CornellClauseRef{14.4}, verify handling an exception against the precise rule category and all cited dependencies.
\end{itemize}\end{CornellReferenceChecklistBox}
\section{Clause Synthesis}
\begin{CornellReferenceOrientationBox}Defines throwing, handler matching, stack unwinding, constructor and destructor behavior, exception specifications, termination, and uncaught-exception observation. Specifies non-local error transfer and cleanup semantics. A sound review distinguishes syntax and participation from runtime preconditions, keeps notes and examples non-mandatory, and follows cross-clause dependencies before making a portability claim.\end{CornellReferenceOrientationBox}
\section{Self-Test Questions}
\begin{enumerate}
\item What role does Preamble play in the clause's overall model?
\item Which normative distinctions must be kept separate when applying Throwing an exception?
\item How would you audit a program or implementation that depends on Constructors and destructors?
\item Compare Handling an exception with Exception specifications; what different question does each answer?
\item What role does Exception specifications play in the clause's overall model?
\item Which normative distinctions must be kept separate when applying Special functions?
\item How would you audit a program or implementation that depends on General?
\item Compare The std::terminate function with The std::uncaught\_\allowbreak{}exceptions function; what different question does each answer?
\item Scenario: a program relies on an unstated implementation behavior while using The std::uncaught\_\allowbreak{}exceptions function. What must be checked before calling the program portable?
\item Scenario: a program relies on an unstated implementation behavior while using exception object. What must be checked before calling the program portable?
\end{enumerate}
\clearpage\section{Answer Key}
\begin{enumerate}
\item Preamble is one part of the clause's core-language semantics model. Its role is understood by connecting the local rule in 14.1 to the clause orientation and its dependent concepts.
\item Keep formation and well-formedness separate from runtime preconditions and effects. Also distinguish mandatory wording from notes, implementation choice from undefined behavior, and interface declarations from detailed contracts; see 14.2.
\item Audit Constructors and destructors by locating the controlling wording in 14.3, listing inputs and type constraints, proving any preconditions, recording effects and failure behavior, and checking lifetime, invalidation, complexity, and synchronization where applicable.
\item Handling an exception and Exception specifications occupy different steps or facility groups. Analyze each under its own subclause, then follow the explicit dependency between them rather than transferring one set of guarantees to the other; begin at 14.4.
\item Exception specifications is one part of the clause's error behavior model. Its role is understood by connecting the local rule in 14.5 to the clause orientation and its dependent concepts.
\item Keep formation and well-formedness separate from runtime preconditions and effects. Also distinguish mandatory wording from notes, implementation choice from undefined behavior, and interface declarations from detailed contracts; see 14.6.
\item Audit General by locating the controlling wording in 14.6.1, listing inputs and type constraints, proving any preconditions, recording effects and failure behavior, and checking lifetime, invalidation, complexity, and synchronization where applicable.
\item The std::terminate function and The std::uncaught\_\allowbreak{}exceptions function occupy different steps or facility groups. Analyze each under its own subclause, then follow the explicit dependency between them rather than transferring one set of guarantees to the other; begin at 14.6.2.
\item First identify the exact requirement in 14.6.3, then classify the relied-on behavior as required, unspecified, implementation-defined, conditionally supported, or undefined. Portability requires satisfying the program obligations and avoiding dependence on a choice not guaranteed in the target environments.
\item First identify the exact requirement in 14, then classify the relied-on behavior as required, unspecified, implementation-defined, conditionally supported, or undefined. Portability requires satisfying the program obligations and avoiding dependence on a choice not guaranteed in the target environments.
\end{enumerate}
\section{Key-Term Recap}
\begin{longtable}{@{}>{\RaggedRight\arraybackslash}p{0.28\textwidth}>{\RaggedRight\arraybackslash}p{0.64\textwidth}@{}}\toprule\textbf{Term} & \textbf{Working meaning}\\\midrule\endfirsthead\toprule\textbf{Term} & \textbf{Working meaning (continued)}\\\midrule\endhead\bottomrule\endfoot
exception object & The object initialized by a throw and inspected by matching handlers.\\
handler & A catch clause selected according to exception matching rules.\\
stack unwinding & Destruction of automatic objects as control exits scopes during exception propagation.\\
exception specification & A function-type property stating whether a function is potentially throwing.\\
termination & The specified failure path used when exception handling cannot continue safely.\\
\end{longtable}
\CornellTakeaway{Exception safety requires proving both the transfer to a valid handler and the lifetime-correct cleanup of every exited scope.}
\end{document}
